% Direct-authority diplomatic French for printed p.206.
\clearpage
\SourcePage{206}

\LayerTarget{tohoku:proof:transformation-first-spectral-sequence:p1}
\noindent de théorème 5.3.1. on a
$\mathbf H^n(G,A_U)(y)=\mathbf H^n(G,A_U(y))$ pour $y\in Y_0$. En vertu de
proposition 5.2.3, l'homomorphisme
$H^n(Y,(A^G)_V)\to H^n(X;G,A_U)$ est donc un isomorphisme pour
tout $n$. On voit tout aussi facilement que $H^n(X;G,A_{X_0})$ s'identifie à
$H^n(X_0;G,A)$, de plus on sait que
$H^n(Y,(A^G)_{Y_0})=H^n(Y_0,A^G)$. On déduit alors du diagramme précédent,
par un ``diagram-chasing'' évident, une suite d'homomorphismes

\LayerTarget{tohoku:sequence:5.4.2}
\[
\tag{5.4.2}
\begin{gathered}
{}\to H^{n-1}(X,G,A)
 \xrightarrow{\beta^{n-1}}
 H^{n-1}(X_0;G,A)/\operatorname{Im}H^{n-1}(Y_0,A^G)
 \xrightarrow{\partial} H^n(Y,A^G)\\[-0.2em]
 \xrightarrow{\alpha^n}H^n(X;G,A)
 \xrightarrow{\beta^n}
 H^n(X_0;G,A)/\operatorname{Im}H^n(Y_0,A^G)
 \to
\end{gathered}
\]
qui est un complexe (le produit de deux homomorphismes consécutifs est nul),
et exacte partout sauf éventuellement aux termes $H^n(Y,A^G)$, le défaut
d'exactitude $\operatorname{Ker}\alpha^n/\operatorname{Im}\partial$ en ce
terme étant isomorphe canoniquement au noyau de l'homomorphisme
$H^n(Y_0,A^G)\to H^n(X_0;G,A)$. On obtient ainsi :

\LayerTarget{tohoku:proposition:5.4.1}
\textsc{Proposition 5.4.1.} \textit{La suite d'homomorphismes \textup{(5.4.2)}
définit sur $H^n(Y,A^G)$ une suite de composition de longueur \textup{3}, dont les
quotients successifs sont isomorphes respectivement à}
$\operatorname{Coker}\beta^{n-1}=H^{n-1}(X_0;G,A)/(\operatorname{Im}H^{n-1}(X;G,A)+\operatorname{Im}H^{n-1}(Y_0,A^G))$,
\textit{à}
$\operatorname{Ker}\bigl(H^n(Y_0,A^G)\to H^n(X_0;G,A)\bigr)$
\textit{et à} $\operatorname{Ker}\beta^n\subset H^n(X;G,A)$.

\LayerTarget{tohoku:comparison:first-and-second-spectral-sequences-5.4}
On peut donc dire que l'utilisation de la première suite spectrale nous donne
un moyen de ramener la détermination de $H^*(Y,A^G)$ à une bonne connaissance
de $H^*(X,G,A)$, $H^*(X_0G,A)$, $H^*(Y_0,A^G)$ et des homomorphismes naturels
$H^*(X;G,A)\to H^*(X_0;G,A)\leftarrow H^*(Y_0,A^G)$.
Or la détermination des deux premiers groupes pourra être tenté par
utilisation de la \textit{deuxième} suite spectrale, tandis que le calcul de
$H^*(Y_0,A^G)$ sera souvent possible si on a pu choisir $Y_0$ assez petit. Un
cas particulièrement simple et important est donné par le

\LayerTarget{tohoku:corollary:5.4.1}
\textsc{Corollaire.} \textit{Supposons que $G$ opère trivialement sur
$X_0$ (donc $G$ est fini), que $A$ soit un faisceau d'espaces vectoriels sur
un corps $k$, et que $G$ opère trivialement sur $A\mathbin{|}X_0$. Alors la
suite \textup{(5.2.4)} est exacte, et $H^n(X_0,G,A)$ s'identifie canoniquement à}
$\displaystyle\bigoplus_{p+q=n}H^p(G,K)\otimes H^q(X_0,A)$.

\LayerTarget{tohoku:proof:corollary-5.4.1}
En effet, la dernière assertion est un cas particulier de la dernière
assertion de théorème 4.4.1. Par suite, l'homomorphisme canonique de
$H^n(Y_0,A^G)=H^n(X_0,A)$ dans $H^n(X_0;G,A)$ est injectif, donc la suite
(5.2.4) est exacte.

\LayerTarget{tohoku:remark:smith-fixed-point-application:p1}
\textsc{Remarques.} Supposons que $G$ soit un groupe fini d'ordre
premier, alors en prenant pour $X_0$ l'ensemble des points de $X$ fixes sous
$G$, et $Y_0=f(X_0)$, $Y_0$ satisfait à la condition énoncée au début de ce
numéro. Si alors $A$ est un faisceau d'espaces vectoriels provenant d'un
faisceau sur $Y$ (i.\ e. tel que $G$ opère trivialement sur
$A\mathbin{|}X_0$) on est dans les conditions d'application du corollaire 1.
On retrouve par exemple très facilement le théorème de P. A. Smith relatif au
cas oû $X$ est une sphère homologique mod. $p$ de dimension finie ($p$ étant
l'ordre de $G$) : alors $X_0$ est une sphère homologique mod. $p$. D'ailleurs
cette démonstration ``naturelle'' par les suites spectrales générales est en
fait très voisine de celle de Borel [2]; bien entendu, l'hypothèse de
compacité faite

\TohokuContinuation{207}
